\scanpage{56}
\mspage{27}
Revenons au cas d'un MD-graphe général $\underline G$. Sa dimension modulaire est
\[
\text{(67)}\qquad\delta(G)=\sum_{\alpha\in S}\underbrace{\bigl(3(g_\alpha-1)+\nu_\alpha+\widetilde\mu_\alpha\bigr)}_{\delta_\alpha}.
\]
Si on veut qu'elle soit égale à $1$ ou $2$, cela signifie que l'on est dans l'un des deux cas :
\[
\left\{\begin{array}{l}
1\text{ ou }2\text{ des }\delta_\alpha\text{ sont égaux à }1,\text{ les autres sont nuls}\\
1\text{ des }\delta_\alpha\text{ est égal à }2,\text{ les autres sont nuls}
\end{array}\right.
\]
Le cas $\delta(G)=1$ correspond au cas où exactement un des $\delta_\alpha$ est $=1$, les autres étant nuls.

II) $\underline G$ \uline{de dim modulaire $1$}

Cela signifie qu'il existe un sommet $\alpha_0$ (et un seul, \deleted{de fait}) de $\underline G$ qui \deleted{soit de genre} \inserted{satisfasse la} condition
\[
\begin{aligned}
\text{(68)}\quad &(\alpha_0\text{ de genre }1,\text{ et son poids total }\nu_{\alpha_0}+\widetilde\mu_{\alpha_0}\text{ est ég. à }\unsure{1},\\
&\qquad\text{i.e. }\alpha_0\text{ est d'ordre }1,\text{ ou sommet isolé de poids marqué }1)\\
&\uline{\text{ou}}\quad(\alpha_0\text{ de genre }0,\text{ et son poids total }\nu_{\alpha_0}+\widetilde\mu_{\alpha_0}=4,\\
&\qquad\text{i.e. }(\nu_{\alpha_0},\widetilde\mu_{\alpha_0})\in\{(4,0),(3,1),(2,2),(1,3),(0,4)\}),
\end{aligned}
\]
avec en plus
\[
\begin{aligned}
\text{(69)}\quad&\text{les sommets }\alpha\ne\alpha_0\text{ sont de genre }0,\\
&\text{d'ordre }\widetilde\mu_\alpha\leq3,\text{ et pour un tel }\alpha\text{ le poids}\\
&\text{marqué }\nu_\alpha\text{ est égal à : }3-\widetilde\mu_\alpha.
\end{aligned}
\]
\deleted{On voit que la MD-structure d'un tel graphe}
\margin{55}
